\begin{helper}
Let \(G\) be an \((n,d,\lambda)\)-loop graph. Then the adjacency action sends the constant-one function to the constant function with value \(d\): for every vertex \(v\),
\[
  (A_G\mathbf 1)(v)=d.
\]
\end{helper}
\begin{proof}
By \(\noderef{LoopGraphNdLambda}\), every vertex of \(G\) has degree \(d\). Fix a vertex \(v\). By \(\noderef{LoopGraphAdjacencyAction}\), \((A_G\mathbf 1)(v)\) is the sum over all vertices \(w\) adjacent from \(v\) of the value \(\mathbf 1(w)=1\), with non-neighbors contributing \(0\). This sum is therefore exactly the number of vertices \(w\) with \(vw\in E(G)\). By \(\noderef{LoopGraphDegree}\), that number is \(\deg_G(v)\), and the degree condition from \(\noderef{LoopGraphNdLambda}\) gives \(\deg_G(v)=d\). Thus \((A_G\mathbf 1)(v)=d\).
\end{proof}
